\begin{center}
\begin{tabular}{lrr}
\toprule
Decisions among 744 hours & Change only$^*$ & Both levels$^*$\\
\midrule
First half strictly better & 40 & 41\\
Second half strictly better & 21 & 21\\
Exact tie & 1 & 1\\
Either strict winner possible & 682 & 681\\
\bottomrule
\end{tabular}\\[3pt]
{\small $^*$Both summaries retain the hourly mean and variation cap.}
\end{center}

Recording the two price levels narrowed the feasible interval in 156 hours. It established just one additional strict winner. No hour in either calculation admitted only a weak ranking with a possible tie. The two-sided witness checks and the direct linear programs agreed on all 1,488 endpoint values to the stated tolerance.

The additional decision occurred in the hour starting at 01:00 market time on 23 January. The original interval in AUD was $[-5261/1200,16909/1200]$. With both levels retained, it collapsed to a single positive value, $104/75$ AUD. The first price was 87.50 and the last 70.86 AUD/MWh. Their variation cap was 44.34 AUD/MWh. Any intervening price $z$ must satisfy $|87.50-z|+|70.86-z|\le44.34$, so $z\ge57.01$. The recorded hourly mean requires the ten interior prices to average exactly 57.01. Each must attain that lower bound. This explains the newly identified winner directly from the market summaries.
